\section*{ID8933: FPGA INTERCONNECT TOPOLOGY Canonical Statutory Formula: G\_FPGA}
\paragraph{Metadata.}
PiID: 3488962766844 | RTEID: RTE:P34688.8228:T0.6078 | UUID: cac70dad-3482-5187-a399-c78f24f8155e

\paragraph{Canonical Statement.}
[ID 928] FPGA INTERCONNECT TOPOLOGY Canonical Statutory Formula: \$G\_\{FPGA\} = Torus\_\{3D\}(10 \textbackslash{}times 10 \textbackslash{}times 1)\$ Formal Technical Dissertation: The graph topology for the control electronics. It connects the FPGA array in a 3D torus structure (nearest neighbor + express lanes) to minimize latency and maximize bandwidth [25]. Source: [25]

\paragraph{Canonical Equation.}
\[
G_{FPGA} = Torus_{3D}(10 \times 10 \times 1)
\]

\paragraph{Dictionary.}
FPGA INTERCONNECT TOPOLOGY Canonical Statutory Formula: G\_FPGA = Torus\_3D(10 × 10 × 1) Formal Technical Dissertation: The graph topology for the control electronics.

\paragraph{Encyclopedia.}
FPGA INTERCONNECT TOPOLOGY Canonical Statutory Formula: G\_FPGA is an algebraic definition, written G\_FPGA = Torus\_3D(10 × 10 × 1). It is an algebraic definition expressing the quantity directly in terms of the variables on its right-hand side. It is built from T, o, r, u, s\_3D, defining G\_FPGA. Within the Trawin Zero Point registry it serves as one element of the framework's mathematical narrative.
